% --- Archon physics formalization source begin ---
% archon:physics
% archon:covers PhyXMiniProblems/problem_phyx_mini_0864.lean
% archon:source-report reports/phyx_mini/problem_phyx_mini_0864.source.json

\chapter{Physics problem phyx\_mini\_0864}
\label{ch:PhyXMiniProblems_problem_phyx_mini_0864}

\paragraph{Problem source.}
\#\# Physical scenario

The four 1.0 g spheres shown in the figure are released simultaneously and allowed to move away from each other.

\#\# Question

What is the speed of each sphere when they are very far apart?

\#\# Answer choices

- A: A: 0.59m/s

- B: B: 1.49m/s

- C: C: 0.19m/s

- D: D: 0.49m/s

\#\# Figure caption (auxiliary; use the image as primary evidence)

The image depicts a schematic diagram of a square with point charges at each corner. 

Content description:

- **Charges**: There are four identical point charges, each labeled as "10 nC" (nanocoulombs), located at the four corners of the square. Each charge is represented by a plus sign within a red circle, indicating they are positive charges.

- **Square Dimensions**: The sides of the square are indicated as "1.0 cm." Each side is represented with dashed lines.

- **Arrangement**: The charges are symmetrically placed, with one at each corner of the square. There is no other text or objects present in the diagram.

The diagram is likely used to illustrate an electrostatic problem involving point charges in a square configuration.

\#\# Dataset metadata

Electrostatics; Physical Model Grounding Reasoning, Multi-Formula Reasoning

\paragraph{Recorded answer/context.}
D: D: 0.49m/s

\paragraph{Figure/image path.}
/root/proposal\_for\_physic/science-mango/phyx\_mini\_run/phyx\_data/test\_image/864.png

\paragraph{Formalization target.}
create a compiling Lean file with sorry bodies at `PhyXMiniProblems/problem\_phyx\_mini\_0864.lean`. The Lean declarations must preserve the physical quantities, dimensions or dimensional roles, figure labels, governing-law hypotheses, and final relation expressed by this problem.
Use Mathlib/Physlib names found through LeanExplore where available. If a physics API is missing, introduce faithful local abstractions rather than scalar placeholder aliases.

\begin{theorem}[Physics formalization target]
\label{thm:physics:phyx_mini_0864:target}
\lean{PhyXMiniProblems.ProblemPhyXMini0864.problem_phyx_mini_0864}
\uses{def:physics:phyx-mini-0864:phyxminiproblems-problemphyxmini0864-speedinmeterspersecond, def:physics:phyx-mini-0864:phyxminiproblems-problemphyxmini0864-spherelabel, def:physics:phyx-mini-0864:phyxminiproblems-problemphyxmini0864-fourchargedspherereleasesetup, def:physics:phyx-mini-0864:phyxminiproblems-problemphyxmini0864-matchesproblemandprimaryfigure, def:physics:phyx-mini-0864:phyxminiproblems-problemphyxmini0864-matchessquaregeometry, def:physics:phyx-mini-0864:phyxminiproblems-problemphyxmini0864-hasphysicalsquarereleaseparameters, def:physics:phyx-mini-0864:phyxminiproblems-problemphyxmini0864-usesvacuumcoulombconstant, def:physics:phyx-mini-0864:phyxminiproblems-problemphyxmini0864-satisfiescoulombandmechanicalenergylaws, def:physics:phyx-mini-0864:phyxminiproblems-problemphyxmini0864-reachesveryfarapartregime, def:physics:phyx-mini-0864:phyxminiproblems-problemphyxmini0864-preservessquarereleasesymmetry, def:physics:phyx-mini-0864:phyxminiproblems-problemphyxmini0864-energyderivedspeedinmeterspersecond, def:physics:phyx-mini-0864:phyxminiproblems-problemphyxmini0864-recordeddatasetanswer, def:physics:phyx-mini-0864:phyxminiproblems-problemphyxmini0864-isuniqueclosestdisplayedspeed}
The assigned autoformalize agent should translate this physics problem into Lean declarations in the covered file, with theorem and lemma proof bodies written as `by sorry`.
\end{theorem}
\begin{proof}
This is an autoformalization task, not a proof task. Produce statements that can later be proved by the physics prover without weakening the physical model.
\end{proof}

% --- Archon declaration topology begin ---
\section{Declaration topology}
\begin{definition}[Length Quantity]
  \label{def:physics:phyx-mini-0864:phyxminiproblems-problemphyxmini0864-lengthquantity}
  \lean{PhyXMiniProblems.ProblemPhyXMini0864.LengthQuantity}
  A nonnegative, unit-independent physical length
\end{definition}
\begin{proof}
  This is the declared model definition of Length Quantity.
\end{proof}
\begin{definition}[Mass Quantity]
  \label{def:physics:phyx-mini-0864:phyxminiproblems-problemphyxmini0864-massquantity}
  \lean{PhyXMiniProblems.ProblemPhyXMini0864.MassQuantity}
  A nonnegative, unit-independent physical mass
\end{definition}
\begin{proof}
  This is the declared model definition of Mass Quantity.
\end{proof}
\begin{definition}[Signed Charge Quantity]
  \label{def:physics:phyx-mini-0864:phyxminiproblems-problemphyxmini0864-signedchargequantity}
  \lean{PhyXMiniProblems.ProblemPhyXMini0864.SignedChargeQuantity}
  A signed, unit-independent electric charge
\end{definition}
\begin{proof}
  This is the declared model definition of Signed Charge Quantity.
\end{proof}
\begin{definition}[speed Dimension]
  \label{def:physics:phyx-mini-0864:phyxminiproblems-problemphyxmini0864-speeddimension}
  \lean{PhyXMiniProblems.ProblemPhyXMini0864.speedDimension}
  The physical dimension L T⁻¹ of speed
\end{definition}
\begin{proof}
  This is the declared model definition of speed Dimension.
\end{proof}
\begin{definition}[Speed Magnitude Quantity]
  \label{def:physics:phyx-mini-0864:phyxminiproblems-problemphyxmini0864-speedmagnitudequantity}
  \lean{PhyXMiniProblems.ProblemPhyXMini0864.SpeedMagnitudeQuantity}
  \uses{def:physics:phyx-mini-0864:phyxminiproblems-problemphyxmini0864-speeddimension}
  A nonnegative physical speed magnitude
\end{definition}
\begin{proof}
  This is the declared model definition of Speed Magnitude Quantity.
\end{proof}
\begin{definition}[length In Meters]
  \label{def:physics:phyx-mini-0864:phyxminiproblems-problemphyxmini0864-lengthinmeters}
  \lean{PhyXMiniProblems.ProblemPhyXMini0864.lengthInMeters}
  \uses{def:physics:phyx-mini-0864:phyxminiproblems-problemphyxmini0864-lengthquantity}
  Coherent-SI metre readout of a physical length
\end{definition}
\begin{proof}
  This is the declared model definition of length In Meters.
\end{proof}
\begin{definition}[length In Centimeters]
  \label{def:physics:phyx-mini-0864:phyxminiproblems-problemphyxmini0864-lengthincentimeters}
  \lean{PhyXMiniProblems.ProblemPhyXMini0864.lengthInCentimeters}
  \uses{def:physics:phyx-mini-0864:phyxminiproblems-problemphyxmini0864-lengthquantity, def:physics:phyx-mini-0864:phyxminiproblems-problemphyxmini0864-lengthinmeters}
  Centimetre readout used by the two dimension labels in the image
\end{definition}
\begin{proof}
  This is the declared model definition of length In Centimeters.
\end{proof}
\begin{definition}[mass In Kilograms]
  \label{def:physics:phyx-mini-0864:phyxminiproblems-problemphyxmini0864-massinkilograms}
  \lean{PhyXMiniProblems.ProblemPhyXMini0864.massInKilograms}
  \uses{def:physics:phyx-mini-0864:phyxminiproblems-problemphyxmini0864-massquantity}
  Coherent-SI kilogram readout of a physical mass
\end{definition}
\begin{proof}
  This is the declared model definition of mass In Kilograms.
\end{proof}
\begin{definition}[mass In Grams]
  \label{def:physics:phyx-mini-0864:phyxminiproblems-problemphyxmini0864-massingrams}
  \lean{PhyXMiniProblems.ProblemPhyXMini0864.massInGrams}
  \uses{def:physics:phyx-mini-0864:phyxminiproblems-problemphyxmini0864-massquantity, def:physics:phyx-mini-0864:phyxminiproblems-problemphyxmini0864-massinkilograms}
  Gram readout used by the prose description of each sphere
\end{definition}
\begin{proof}
  This is the declared model definition of mass In Grams.
\end{proof}
\begin{definition}[charge In Coulombs]
  \label{def:physics:phyx-mini-0864:phyxminiproblems-problemphyxmini0864-chargeincoulombs}
  \lean{PhyXMiniProblems.ProblemPhyXMini0864.chargeInCoulombs}
  \uses{def:physics:phyx-mini-0864:phyxminiproblems-problemphyxmini0864-signedchargequantity}
  Coherent-SI coulomb readout of a signed physical charge
\end{definition}
\begin{proof}
  This is the declared model definition of charge In Coulombs.
\end{proof}
\begin{definition}[charge In Nanocoulombs]
  \label{def:physics:phyx-mini-0864:phyxminiproblems-problemphyxmini0864-chargeinnanocoulombs}
  \lean{PhyXMiniProblems.ProblemPhyXMini0864.chargeInNanocoulombs}
  \uses{def:physics:phyx-mini-0864:phyxminiproblems-problemphyxmini0864-signedchargequantity, def:physics:phyx-mini-0864:phyxminiproblems-problemphyxmini0864-chargeincoulombs}
  Nanocoulomb readout used by the four printed charge labels
\end{definition}
\begin{proof}
  This is the declared model definition of charge In Nanocoulombs.
\end{proof}
\begin{definition}[speed In Meters Per Second]
  \label{def:physics:phyx-mini-0864:phyxminiproblems-problemphyxmini0864-speedinmeterspersecond}
  \lean{PhyXMiniProblems.ProblemPhyXMini0864.speedInMetersPerSecond}
  \uses{def:physics:phyx-mini-0864:phyxminiproblems-problemphyxmini0864-speedmagnitudequantity}
  Coherent-SI metre-per-second readout of a physical speed
\end{definition}
\begin{proof}
  This is the declared model definition of speed In Meters Per Second.
\end{proof}
\begin{definition}[energy In Joules]
  \label{def:physics:phyx-mini-0864:phyxminiproblems-problemphyxmini0864-energyinjoules}
  \lean{PhyXMiniProblems.ProblemPhyXMini0864.energyInJoules}
  Coherent-SI joule readout of a physical energy
\end{definition}
\begin{proof}
  This is the declared model definition of energy In Joules.
\end{proof}
\begin{definition}[Sphere Label]
  \label{def:physics:phyx-mini-0864:phyxminiproblems-problemphyxmini0864-spherelabel}
  \lean{PhyXMiniProblems.ProblemPhyXMini0864.SphereLabel}
  The four physical spheres, named by their initial image locations
\end{definition}
\begin{proof}
  This is the declared model definition of Sphere Label.
\end{proof}
\begin{definition}[Square Corner]
  \label{def:physics:phyx-mini-0864:phyxminiproblems-problemphyxmini0864-squarecorner}
  \lean{PhyXMiniProblems.ProblemPhyXMini0864.SquareCorner}
  The four corners of the dashed square
\end{definition}
\begin{proof}
  This is the declared model definition of Square Corner.
\end{proof}
\begin{definition}[expected Sphere Corner]
  \label{def:physics:phyx-mini-0864:phyxminiproblems-problemphyxmini0864-expectedspherecorner}
  \lean{PhyXMiniProblems.ProblemPhyXMini0864.expectedSphereCorner}
  \uses{def:physics:phyx-mini-0864:phyxminiproblems-problemphyxmini0864-spherelabel, def:physics:phyx-mini-0864:phyxminiproblems-problemphyxmini0864-squarecorner}
  Corner at which each named sphere appears in the primary image
\end{definition}
\begin{proof}
  This is the declared model definition of expected Sphere Corner.
\end{proof}
\begin{definition}[Sphere Pair]
  \label{def:physics:phyx-mini-0864:phyxminiproblems-problemphyxmini0864-spherepair}
  \lean{PhyXMiniProblems.ProblemPhyXMini0864.SpherePair}
  The six unordered pairs among the four spheres
\end{definition}
\begin{proof}
  This is the declared model definition of Sphere Pair.
\end{proof}
\begin{definition}[first]
  \label{def:physics:phyx-mini-0864:phyxminiproblems-problemphyxmini0864-spherepair-first}
  \lean{PhyXMiniProblems.ProblemPhyXMini0864.SpherePair.first}
  \uses{def:physics:phyx-mini-0864:phyxminiproblems-problemphyxmini0864-spherelabel, def:physics:phyx-mini-0864:phyxminiproblems-problemphyxmini0864-spherepair}
  First endpoint of a named unordered sphere pair
\end{definition}
\begin{proof}
  This is the declared model definition of first.
\end{proof}
\begin{definition}[second]
  \label{def:physics:phyx-mini-0864:phyxminiproblems-problemphyxmini0864-spherepair-second}
  \lean{PhyXMiniProblems.ProblemPhyXMini0864.SpherePair.second}
  \uses{def:physics:phyx-mini-0864:phyxminiproblems-problemphyxmini0864-spherelabel, def:physics:phyx-mini-0864:phyxminiproblems-problemphyxmini0864-spherepair}
  Second endpoint of a named unordered sphere pair
\end{definition}
\begin{proof}
  This is the declared model definition of second.
\end{proof}
\begin{definition}[Pair Geometry]
  \label{def:physics:phyx-mini-0864:phyxminiproblems-problemphyxmini0864-pairgeometry}
  \lean{PhyXMiniProblems.ProblemPhyXMini0864.PairGeometry}
  Whether a pair spans a side or a diagonal of the square
\end{definition}
\begin{proof}
  This is the declared model definition of Pair Geometry.
\end{proof}
\begin{definition}[geometry]
  \label{def:physics:phyx-mini-0864:phyxminiproblems-problemphyxmini0864-spherepair-geometry}
  \lean{PhyXMiniProblems.ProblemPhyXMini0864.SpherePair.geometry}
  \uses{def:physics:phyx-mini-0864:phyxminiproblems-problemphyxmini0864-spherepair, def:physics:phyx-mini-0864:phyxminiproblems-problemphyxmini0864-pairgeometry}
  Geometric role of each unordered sphere pair
\end{definition}
\begin{proof}
  This is the declared model definition of geometry.
\end{proof}
\begin{definition}[Square Edge]
  \label{def:physics:phyx-mini-0864:phyxminiproblems-problemphyxmini0864-squareedge}
  \lean{PhyXMiniProblems.ProblemPhyXMini0864.SquareEdge}
  The four dashed boundary edges visible in the raster
\end{definition}
\begin{proof}
  This is the declared model definition of Square Edge.
\end{proof}
\begin{definition}[expected Dimension Label Centimeters]
  \label{def:physics:phyx-mini-0864:phyxminiproblems-problemphyxmini0864-expecteddimensionlabelcentimeters}
  \lean{PhyXMiniProblems.ProblemPhyXMini0864.expectedDimensionLabelCentimeters}
  \uses{def:physics:phyx-mini-0864:phyxminiproblems-problemphyxmini0864-squareedge}
  The primary image prints 1.0 cm only on the top and left edges
\end{definition}
\begin{proof}
  This is the declared model definition of expected Dimension Label Centimeters.
\end{proof}
\begin{definition}[Charge Circle Color]
  \label{def:physics:phyx-mini-0864:phyxminiproblems-problemphyxmini0864-chargecirclecolor}
  \lean{PhyXMiniProblems.ProblemPhyXMini0864.ChargeCircleColor}
  Fill colour of each positive-charge circle in the supplied image
\end{definition}
\begin{proof}
  This is the declared model definition of Charge Circle Color.
\end{proof}
\begin{definition}[Charge Sign Glyph]
  \label{def:physics:phyx-mini-0864:phyxminiproblems-problemphyxmini0864-chargesignglyph}
  \lean{PhyXMiniProblems.ProblemPhyXMini0864.ChargeSignGlyph}
  Charge-sign glyph drawn inside every sphere circle
\end{definition}
\begin{proof}
  This is the declared model definition of Charge Sign Glyph.
\end{proof}
\begin{definition}[Four Charge Square Figure]
  \label{def:physics:phyx-mini-0864:phyxminiproblems-problemphyxmini0864-fourchargesquarefigure}
  \lean{PhyXMiniProblems.ProblemPhyXMini0864.FourChargeSquareFigure}
  \uses{def:physics:phyx-mini-0864:phyxminiproblems-problemphyxmini0864-spherelabel, def:physics:phyx-mini-0864:phyxminiproblems-problemphyxmini0864-squarecorner, def:physics:phyx-mini-0864:phyxminiproblems-problemphyxmini0864-squareedge, def:physics:phyx-mini-0864:phyxminiproblems-problemphyxmini0864-chargecirclecolor, def:physics:phyx-mini-0864:phyxminiproblems-problemphyxmini0864-chargesignglyph}
  Literal presentation data transcribed from image 864.png
\end{definition}
\begin{proof}
  This is the declared model definition of Four Charge Square Figure.
\end{proof}
\begin{definition}[Interaction Model]
  \label{def:physics:phyx-mini-0864:phyxminiproblems-problemphyxmini0864-interactionmodel}
  \lean{PhyXMiniProblems.ProblemPhyXMini0864.InteractionModel}
  The qualitative interaction idealization used after simultaneous release
\end{definition}
\begin{proof}
  This is the declared model definition of Interaction Model.
\end{proof}
\begin{definition}[Separation Regime]
  \label{def:physics:phyx-mini-0864:phyxminiproblems-problemphyxmini0864-separationregime}
  \lean{PhyXMiniProblems.ProblemPhyXMini0864.SeparationRegime}
  The separation regime at which the final speed is observed
\end{definition}
\begin{proof}
  This is the declared model definition of Separation Regime.
\end{proof}
\begin{definition}[Four Charged Sphere Release Setup]
  \label{def:physics:phyx-mini-0864:phyxminiproblems-problemphyxmini0864-fourchargedspherereleasesetup}
  \lean{PhyXMiniProblems.ProblemPhyXMini0864.FourChargedSphereReleaseSetup}
  \uses{def:physics:phyx-mini-0864:phyxminiproblems-problemphyxmini0864-lengthquantity, def:physics:phyx-mini-0864:phyxminiproblems-problemphyxmini0864-massquantity, def:physics:phyx-mini-0864:phyxminiproblems-problemphyxmini0864-signedchargequantity, def:physics:phyx-mini-0864:phyxminiproblems-problemphyxmini0864-speedmagnitudequantity, def:physics:phyx-mini-0864:phyxminiproblems-problemphyxmini0864-spherelabel, def:physics:phyx-mini-0864:phyxminiproblems-problemphyxmini0864-spherepair, def:physics:phyx-mini-0864:phyxminiproblems-problemphyxmini0864-fourchargesquarefigure, def:physics:phyx-mini-0864:phyxminiproblems-problemphyxmini0864-interactionmodel, def:physics:phyx-mini-0864:phyxminiproblems-problemphyxmini0864-separationregime}
  Independent quantities and energy observables for the square release. The final speeds are unknown physical observables. They are not defined from the recorded answer or from the derived square-root expression
\end{definition}
\begin{proof}
  This is the declared model definition of Four Charged Sphere Release Setup.
\end{proof}
\begin{definition}[position Vector In Meters]
  \label{def:physics:phyx-mini-0864:phyxminiproblems-problemphyxmini0864-positionvectorinmeters}
  \lean{PhyXMiniProblems.ProblemPhyXMini0864.positionVectorInMeters}
  Turn a planar Space 2 point into its explicitly metre-valued vector
\end{definition}
\begin{proof}
  This is the declared model definition of position Vector In Meters.
\end{proof}
\begin{definition}[initial Pair Displacement In Meters]
  \label{def:physics:phyx-mini-0864:phyxminiproblems-problemphyxmini0864-initialpairdisplacementinmeters}
  \lean{PhyXMiniProblems.ProblemPhyXMini0864.initialPairDisplacementInMeters}
  \uses{def:physics:phyx-mini-0864:phyxminiproblems-problemphyxmini0864-spherepair, def:physics:phyx-mini-0864:phyxminiproblems-problemphyxmini0864-fourchargedspherereleasesetup, def:physics:phyx-mini-0864:phyxminiproblems-problemphyxmini0864-positionvectorinmeters}
  Initial displacement between the endpoints of an unordered pair
\end{definition}
\begin{proof}
  This is the declared model definition of initial Pair Displacement In Meters.
\end{proof}
\begin{definition}[Matches Problem And Primary Figure]
  \label{def:physics:phyx-mini-0864:phyxminiproblems-problemphyxmini0864-matchesproblemandprimaryfigure}
  \lean{PhyXMiniProblems.ProblemPhyXMini0864.MatchesProblemAndPrimaryFigure}
  \uses{def:physics:phyx-mini-0864:phyxminiproblems-problemphyxmini0864-lengthincentimeters, def:physics:phyx-mini-0864:phyxminiproblems-problemphyxmini0864-massingrams, def:physics:phyx-mini-0864:phyxminiproblems-problemphyxmini0864-chargeinnanocoulombs, def:physics:phyx-mini-0864:phyxminiproblems-problemphyxmini0864-speedinmeterspersecond, def:physics:phyx-mini-0864:phyxminiproblems-problemphyxmini0864-expectedspherecorner, def:physics:phyx-mini-0864:phyxminiproblems-problemphyxmini0864-expecteddimensionlabelcentimeters, def:physics:phyx-mini-0864:phyxminiproblems-problemphyxmini0864-fourchargedspherereleasesetup}
  The numerical labels, corner associations, and release conditions stated in the problem and visible in the primary image. No final speed occurs here
\end{definition}
\begin{proof}
  This is the declared model definition of Matches Problem And Primary Figure.
\end{proof}
\begin{definition}[Matches Square Geometry]
  \label{def:physics:phyx-mini-0864:phyxminiproblems-problemphyxmini0864-matchessquaregeometry}
  \lean{PhyXMiniProblems.ProblemPhyXMini0864.MatchesSquareGeometry}
  \uses{def:physics:phyx-mini-0864:phyxminiproblems-problemphyxmini0864-lengthinmeters, def:physics:phyx-mini-0864:phyxminiproblems-problemphyxmini0864-fourchargedspherereleasesetup, def:physics:phyx-mini-0864:phyxminiproblems-problemphyxmini0864-positionvectorinmeters, def:physics:phyx-mini-0864:phyxminiproblems-problemphyxmini0864-initialpairdisplacementinmeters}
  Coordinate realization of the square and the relation between Euclidean separation and the physical pair-distance quantity. The bottom-left corner is chosen as the coordinate origin without loss of physical generality
\end{definition}
\begin{proof}
  This is the declared model definition of Matches Square Geometry.
\end{proof}
\begin{definition}[Has Physical Square Release Parameters]
  \label{def:physics:phyx-mini-0864:phyxminiproblems-problemphyxmini0864-hasphysicalsquarereleaseparameters}
  \lean{PhyXMiniProblems.ProblemPhyXMini0864.HasPhysicalSquareReleaseParameters}
  \uses{def:physics:phyx-mini-0864:phyxminiproblems-problemphyxmini0864-lengthinmeters, def:physics:phyx-mini-0864:phyxminiproblems-problemphyxmini0864-massinkilograms, def:physics:phyx-mini-0864:phyxminiproblems-problemphyxmini0864-chargeincoulombs, def:physics:phyx-mini-0864:phyxminiproblems-problemphyxmini0864-fourchargedspherereleasesetup}
  Positivity conditions selecting the nondegenerate physical branch
\end{definition}
\begin{proof}
  This is the declared model definition of Has Physical Square Release Parameters.
\end{proof}
\begin{definition}[Uses Vacuum Coulomb Constant]
  \label{def:physics:phyx-mini-0864:phyxminiproblems-problemphyxmini0864-usesvacuumcoulombconstant}
  \lean{PhyXMiniProblems.ProblemPhyXMini0864.UsesVacuumCoulombConstant}
  \uses{def:physics:phyx-mini-0864:phyxminiproblems-problemphyxmini0864-fourchargedspherereleasesetup}
  Free-space calibration of Physlib's electromagnetic system in SI units
\end{definition}
\begin{proof}
  This is the declared model definition of Uses Vacuum Coulomb Constant.
\end{proof}
\begin{definition}[Satisfies Coulomb And Mechanical Energy Laws]
  \label{def:physics:phyx-mini-0864:phyxminiproblems-problemphyxmini0864-satisfiescoulombandmechanicalenergylaws}
  \lean{PhyXMiniProblems.ProblemPhyXMini0864.SatisfiesCoulombAndMechanicalEnergyLaws}
  \uses{def:physics:phyx-mini-0864:phyxminiproblems-problemphyxmini0864-lengthinmeters, def:physics:phyx-mini-0864:phyxminiproblems-problemphyxmini0864-massinkilograms, def:physics:phyx-mini-0864:phyxminiproblems-problemphyxmini0864-chargeincoulombs, def:physics:phyx-mini-0864:phyxminiproblems-problemphyxmini0864-speedinmeterspersecond, def:physics:phyx-mini-0864:phyxminiproblems-problemphyxmini0864-energyinjoules, def:physics:phyx-mini-0864:phyxminiproblems-problemphyxmini0864-spherelabel, def:physics:phyx-mini-0864:phyxminiproblems-problemphyxmini0864-spherepair, def:physics:phyx-mini-0864:phyxminiproblems-problemphyxmini0864-fourchargedspherereleasesetup}
  Every unordered pair contributes k q₁ q₂/r, the pair energies add to the initial electrostatic energy, and each sphere has translational kinetic energy m v²/2. Conservation is stated between release and the far-apart state. These laws contain no numerical assertion about the final speed
\end{definition}
\begin{proof}
  This is the declared model definition of Satisfies Coulomb And Mechanical Energy Laws.
\end{proof}
\begin{definition}[Reaches Very Far Apart Regime]
  \label{def:physics:phyx-mini-0864:phyxminiproblems-problemphyxmini0864-reachesveryfarapartregime}
  \lean{PhyXMiniProblems.ProblemPhyXMini0864.ReachesVeryFarApartRegime}
  \uses{def:physics:phyx-mini-0864:phyxminiproblems-problemphyxmini0864-energyinjoules, def:physics:phyx-mini-0864:phyxminiproblems-problemphyxmini0864-fourchargedspherereleasesetup}
  At asymptotically large mutual separation the zero of Coulomb potential is chosen so that the remaining electrostatic potential energy vanishes
\end{definition}
\begin{proof}
  This is the declared model definition of Reaches Very Far Apart Regime.
\end{proof}
\begin{definition}[Preserves Square Release Symmetry]
  \label{def:physics:phyx-mini-0864:phyxminiproblems-problemphyxmini0864-preservessquarereleasesymmetry}
  \lean{PhyXMiniProblems.ProblemPhyXMini0864.PreservesSquareReleaseSymmetry}
  \uses{def:physics:phyx-mini-0864:phyxminiproblems-problemphyxmini0864-speedinmeterspersecond, def:physics:phyx-mini-0864:phyxminiproblems-problemphyxmini0864-fourchargedspherereleasesetup}
  The isolated evolution preserves the rotational and reflection symmetries of four identical charges released from rest at the corners of a square. This states equality of final speed magnitudes, not their requested value
\end{definition}
\begin{proof}
  This is the declared model definition of Preserves Square Release Symmetry.
\end{proof}
\begin{definition}[energy Derived Speed In Meters Per Second]
  \label{def:physics:phyx-mini-0864:phyxminiproblems-problemphyxmini0864-energyderivedspeedinmeterspersecond}
  \lean{PhyXMiniProblems.ProblemPhyXMini0864.energyDerivedSpeedInMetersPerSecond}
  \uses{def:physics:phyx-mini-0864:phyxminiproblems-problemphyxmini0864-lengthinmeters, def:physics:phyx-mini-0864:phyxminiproblems-problemphyxmini0864-massinkilograms, def:physics:phyx-mini-0864:phyxminiproblems-problemphyxmini0864-chargeincoulombs, def:physics:phyx-mini-0864:phyxminiproblems-problemphyxmini0864-fourchargedspherereleasesetup}
  The speed candidate obtained by converting the six initial pair energies into the kinetic energy of four equal-mass spheres. The factor 4 + sqrt 2 comes from four side pairs and two diagonal pairs
\end{definition}
\begin{proof}
  This is the declared model definition of energy Derived Speed In Meters Per Second.
\end{proof}
\begin{definition}[Answer Choice]
  \label{def:physics:phyx-mini-0864:phyxminiproblems-problemphyxmini0864-answerchoice}
  \lean{PhyXMiniProblems.ProblemPhyXMini0864.AnswerChoice}
  Labels of the four speed choices printed with the problem
\end{definition}
\begin{proof}
  This is the declared model definition of Answer Choice.
\end{proof}
\begin{definition}[speed In Meters Per Second]
  \label{def:physics:phyx-mini-0864:phyxminiproblems-problemphyxmini0864-answerchoice-speedinmeterspersecond}
  \lean{PhyXMiniProblems.ProblemPhyXMini0864.AnswerChoice.speedInMetersPerSecond}
  \uses{def:physics:phyx-mini-0864:phyxminiproblems-problemphyxmini0864-speedinmeterspersecond, def:physics:phyx-mini-0864:phyxminiproblems-problemphyxmini0864-answerchoice}
  Metre-per-second value printed beside each answer label
\end{definition}
\begin{proof}
  This is the declared model definition of speed In Meters Per Second.
\end{proof}
\begin{definition}[recorded Dataset Answer]
  \label{def:physics:phyx-mini-0864:phyxminiproblems-problemphyxmini0864-recordeddatasetanswer}
  \lean{PhyXMiniProblems.ProblemPhyXMini0864.recordedDatasetAnswer}
  \uses{def:physics:phyx-mini-0864:phyxminiproblems-problemphyxmini0864-answerchoice}
  The answer label recorded by the source dataset
\end{definition}
\begin{proof}
  This is the declared model definition of recorded Dataset Answer.
\end{proof}
\begin{definition}[Is Unique Closest Displayed Speed]
  \label{def:physics:phyx-mini-0864:phyxminiproblems-problemphyxmini0864-isuniqueclosestdisplayedspeed}
  \lean{PhyXMiniProblems.ProblemPhyXMini0864.IsUniqueClosestDisplayedSpeed}
  \uses{def:physics:phyx-mini-0864:phyxminiproblems-problemphyxmini0864-speedinmeterspersecond, def:physics:phyx-mini-0864:phyxminiproblems-problemphyxmini0864-spherelabel, def:physics:phyx-mini-0864:phyxminiproblems-problemphyxmini0864-fourchargedspherereleasesetup, def:physics:phyx-mini-0864:phyxminiproblems-problemphyxmini0864-answerchoice}
  A displayed answer is uniquely closest to every sphere's final speed
\end{definition}
\begin{proof}
  This is the declared model definition of Is Unique Closest Displayed Speed.
\end{proof}
\begin{lemma}[initial Pair Separation formula]
  \label{lem:physics:phyx-mini-0864:phyxminiproblems-problemphyxmini0864-initialpairseparation-formula}
  \lean{PhyXMiniProblems.ProblemPhyXMini0864.initialPairSeparation_formula}
  \uses{def:physics:phyx-mini-0864:phyxminiproblems-problemphyxmini0864-lengthinmeters, def:physics:phyx-mini-0864:phyxminiproblems-problemphyxmini0864-spherepair, def:physics:phyx-mini-0864:phyxminiproblems-problemphyxmini0864-fourchargedspherereleasesetup, def:physics:phyx-mini-0864:phyxminiproblems-problemphyxmini0864-matchessquaregeometry}
  Euclidean square geometry gives four separations equal to the side length and two equal to sqrt 2 times that length. This is a derived geometric fact, not an extra setup assumption
\end{lemma}
\begin{proof}
  The conclusion follows from the governing relations and figure readouts named in the statement of initial Pair Separation formula.
\end{proof}
\begin{lemma}[initial Electrostatic Potential Energy formula]
  \label{lem:physics:phyx-mini-0864:phyxminiproblems-problemphyxmini0864-initialelectrostaticpotentialenergy-formula}
  \lean{PhyXMiniProblems.ProblemPhyXMini0864.initialElectrostaticPotentialEnergy_formula}
  \uses{def:physics:phyx-mini-0864:phyxminiproblems-problemphyxmini0864-lengthinmeters, def:physics:phyx-mini-0864:phyxminiproblems-problemphyxmini0864-chargeincoulombs, def:physics:phyx-mini-0864:phyxminiproblems-problemphyxmini0864-energyinjoules, def:physics:phyx-mini-0864:phyxminiproblems-problemphyxmini0864-fourchargedspherereleasesetup, def:physics:phyx-mini-0864:phyxminiproblems-problemphyxmini0864-matchesproblemandprimaryfigure, def:physics:phyx-mini-0864:phyxminiproblems-problemphyxmini0864-matchessquaregeometry, def:physics:phyx-mini-0864:phyxminiproblems-problemphyxmini0864-hasphysicalsquarereleaseparameters, def:physics:phyx-mini-0864:phyxminiproblems-problemphyxmini0864-satisfiescoulombandmechanicalenergylaws}
  Summing the four edge-pair and two diagonal-pair Coulomb energies yields the initial electrostatic energy. No displayed speed is used
\end{lemma}
\begin{proof}
  The conclusion follows from the governing relations and figure readouts named in the statement of initial Electrostatic Potential Energy formula.
\end{proof}
\begin{lemma}[final Speed formula]
  \label{lem:physics:phyx-mini-0864:phyxminiproblems-problemphyxmini0864-finalspeed-formula}
  \lean{PhyXMiniProblems.ProblemPhyXMini0864.finalSpeed_formula}
  \uses{def:physics:phyx-mini-0864:phyxminiproblems-problemphyxmini0864-speedinmeterspersecond, def:physics:phyx-mini-0864:phyxminiproblems-problemphyxmini0864-spherelabel, def:physics:phyx-mini-0864:phyxminiproblems-problemphyxmini0864-fourchargedspherereleasesetup, def:physics:phyx-mini-0864:phyxminiproblems-problemphyxmini0864-matchesproblemandprimaryfigure, def:physics:phyx-mini-0864:phyxminiproblems-problemphyxmini0864-matchessquaregeometry, def:physics:phyx-mini-0864:phyxminiproblems-problemphyxmini0864-hasphysicalsquarereleaseparameters, def:physics:phyx-mini-0864:phyxminiproblems-problemphyxmini0864-satisfiescoulombandmechanicalenergylaws, def:physics:phyx-mini-0864:phyxminiproblems-problemphyxmini0864-reachesveryfarapartregime, def:physics:phyx-mini-0864:phyxminiproblems-problemphyxmini0864-preservessquarereleasesymmetry, def:physics:phyx-mini-0864:phyxminiproblems-problemphyxmini0864-energyderivedspeedinmeterspersecond}
  Energy conservation, release from rest, the far-apart zero of potential, and square symmetry determine the common final speed
\end{lemma}
\begin{proof}
  The conclusion follows from the governing relations and figure readouts named in the statement of final Speed formula.
\end{proof}
\begin{lemma}[final Speed numerical Bounds]
  \label{lem:physics:phyx-mini-0864:phyxminiproblems-problemphyxmini0864-finalspeed-numericalbounds}
  \lean{PhyXMiniProblems.ProblemPhyXMini0864.finalSpeed_numericalBounds}
  \uses{def:physics:phyx-mini-0864:phyxminiproblems-problemphyxmini0864-speedinmeterspersecond, def:physics:phyx-mini-0864:phyxminiproblems-problemphyxmini0864-spherelabel, def:physics:phyx-mini-0864:phyxminiproblems-problemphyxmini0864-fourchargedspherereleasesetup, def:physics:phyx-mini-0864:phyxminiproblems-problemphyxmini0864-matchesproblemandprimaryfigure, def:physics:phyx-mini-0864:phyxminiproblems-problemphyxmini0864-matchessquaregeometry, def:physics:phyx-mini-0864:phyxminiproblems-problemphyxmini0864-hasphysicalsquarereleaseparameters, def:physics:phyx-mini-0864:phyxminiproblems-problemphyxmini0864-usesvacuumcoulombconstant, def:physics:phyx-mini-0864:phyxminiproblems-problemphyxmini0864-satisfiescoulombandmechanicalenergylaws, def:physics:phyx-mini-0864:phyxminiproblems-problemphyxmini0864-reachesveryfarapartregime, def:physics:phyx-mini-0864:phyxminiproblems-problemphyxmini0864-preservessquarereleasesymmetry}
  With the printed mass, charge, side length, and reference Coulomb constant, the derived final speed lies strictly between 0.49 m/s and 0.50 m/s
\end{lemma}
\begin{proof}
  The conclusion follows from the governing relations and figure readouts named in the statement of final Speed numerical Bounds.
\end{proof}
% --- Archon declaration topology end ---
% --- Archon physics formalization source end ---
